 %%%%%%%%%%%%%%%%%%%%%%%%%%%%%%%%%%%%%%%%%%%%%%%%%%%%%%%%%%%%%%%%%
%%% %
%%% % weiiszablon.tex
%%% % The Faculty of Electrical and Computer Engineering
%%% % Rzeszow University Of Technology diploma thesis Template
%%% % Szablon pracy dyplomowej Wydziału Elektrotechniki 
%%% % i Informatyki PRz
%%% % January, 2024
%%%%%%%%%%%%%%%%%%%%%%%%%%%%%%%%%%%%%%%%%%%%%%%%%%%%%%%%%%%%%%%%%

\documentclass[12pt,twoside]{article}

\usepackage{weiiszablon}
\usepackage{tabularx}
\usepackage{booktabs}
\usepackage{indentfirst}

\author{Jakub Portas, Mateusz Skali, Katarzyna Materia, Magdalena Matuła, Wiktor Kuczek}

% np. EF-123456, EN-654321, ...
\studentID{}

\title{Porównanie platform sprzętowych w detekcji anomalii dla przykładowych zbiorów - Jetson, RPI, ESP32, PC}
\titleEN{}


%%% wybierz rodzaj pracy wpisując jeden z poniższych numerów: ...
% 1 = inżynierska	% BSc
% 2 = magisterska	% MSc
% 3 = doktorska		% PhD
% 4 = praca inżynierska
%%% na miejsce zera w linijce poniżej
\newcommand{\rodzajPracyNo}{5}



\abstract{Treść streszczenia po polsku}
\abstractEN{Treść streszczenia po angielsku}

\keywords{max. 5 słów kluczowych w 2 wierszach, oddzielanych przecinkami}
\keywordsEN{max. five keywords in English}


\begin{document}

% strona tytułowa
\maketitle

% spis treści
\tableofcontents

\clearpage


\section{Wstęp i cel projektu}
W dobie czwartej rewolucji przemysłowej oraz gwałtownego rozwoju systemów Internetu Rzeczy, niezawodne monitorowanie stanu maszyn i wczesne wykrywanie usterek stały się kluczowymi wyzwaniami inżynieryjnymi. Tradycyjne podejście do analizy danych, polegające na przesyłaniu ogromnych ilości surowych odczytów z sensorów do scentralizowanej chmury obliczeniowej, generuje opóźnienia, problemy z przepustowością sieci oraz ryzyka związane z bezpieczeństwem danych \cite{shi2016edge}. Z tego powodu coraz większą popularność zyskuje koncepcja przetwarzania brzegowego (Edge Computing), która zakłada przeniesienie algorytmów uczenia maszynowego bezpośrednio na urządzenia zlokalizowane blisko źródła sygnału. Wymaga to jednak starannego doboru metod analitycznych do mocno ograniczonych zasobów sprzętowych tych urządzeń.

Wdrożenie przetwarzania brzegowego napotyka jednak na istotny problem inżynieryjny, czyli ogromną fragmentację rynku platform sprzętowych. Współczesny ekosystem Edge obejmuje diametralnie różne architektury od energooszczędnych mikrokontrolerów z niewielką ilością pamięci, przez komputery jednopłytkowe (SBC) ogólnego przeznaczenia oparte na architekturze ARM, aż po wysoce wyspecjalizowane, przemysłowe moduły wyposażone w akceleratory sztucznej inteligencji. Przeniesienie oprogramowania analitycznego ze środowiska stacji roboczej na te urządzenia nierzadko wiąże się ze spadkiem wydajności, problemami z architekturą oprogramowania czy brakiem sprzętowego wsparcia dla określonych instrukcji.

Głównym celem niniejszego projektu grupowego jest wieloaspektowa ewaluacja i porównanie wydajności czterech zróżnicowanych platform sprzętowych w kontekście ich przydatności do zadań przemysłowego uczenia maszynowego. Do testów wytypowano: klasyczny komputer osobisty z procesorem architektury x86\_64, nowoczesny komputer jednopłytkowy Raspberry Pi 5, dedykowany moduł Edge w postaci NVIDIA Jetson Orin Nano oraz mikrokontroler z rodziny ESP32. W celu rzetelnego zmierzenia ich możliwości, jako znormalizowane obciążenie testowe (tzw. benchmark) wykorzystano zestaw trzech klasycznych algorytmów detekcji anomalii: Autoenkoder, Isolation Forest oraz strumieniowy model Half-Space Trees. Modele te, różniące się złożonością czasową i zapotrzebowaniem na pamięć operacyjną, posłużyły jako miara do obiektywnego obciążenia badanych architektur procesorów.

Zakres prac koncentruje się przede wszystkim na rygorystycznych testach wydajnościowych sprzętu. Ocenie poddano czas wykonania pełnego potoku danych, opóźnienia (latencję) w przetwarzaniu pojedynczych próbek strumienia oraz maksymalny narzut na zasoby systemowe maszyny. Pomiary jakości samej klasyfikacji pełnią w tym projekcie funkcję walidacyjną. Mają za zadanie udowodnić stabilność i matematyczną tożsamość modeli po ich skompilowaniu na różnych architekturach sprzętowych (ARM vs x86). Kluczowym elementem pracy jest również analiza trudności środowiskowych, takich jak ograniczenia pamięci zunifikowanej, blokady systemowe czy kompatybilność bibliotek języka Python z akceleracją sprzętową. Wyniki tych badań pozwolą na sformułowanie praktycznych wniosków dotyczących doboru optymalnej platformy brzegowej do konkretnych wymagań systemów detekcji anomalii.

\section{Charakterystyka środowisk testowych}
Aby w pełni zbadać wpływ architektury sprzętowej na wydajność przetwarzania danych, do projektu wytypowano urządzenia reprezentujące cztery skrajnie odmienne środowiska obliczeniowe. Zestawienie to pokrywa pełne spektrum współczesnej inżynierii systemów od scentralizowanych stacji roboczych x86\_64, przez uniwersalne komputery jednopłytkowe i wyspecjalizowane moduły akcelerowane sprzętowo, aż po skrajnie ograniczone zasobowo mikrokontrolery. Badane algorytmy posłużyły na każdej z tych platform jako znormalizowane obciążenie testowe, mające na celu weryfikację ich limitów operacyjnych.
\subsection{Stacja robocza x86\_64 (PC) – Środowisko referencyjne}
Jako sprzętowy punkt odniesienia wykorzystano wysokowydajny komputer osobisty Lenovo LOQ. Architektura ta reprezentuje w projekcie klasyczne serwery chmurowe. Wykorzystanie procesora w architekturze x86\_64, zasilania sieciowego oraz aktywnego układu chłodzenia eliminuje wszelkie ograniczenia związane z dławieniem termicznym oraz oszczędzaniem energii. Głównym celem implementacji algorytmów na tym urządzeniu było wyznaczenie maksymalnej, teoretycznej przepustowości potoku danych dla użytych bibliotek w języku Python. Wyniki z PC stanowią "złoty standard", pozwalający na precyzyjne zmierzenie, jak duży narzut wydajnościowy w postaci opóźnień lub spadku liczby operacji na sekundę pociąga za sobą przeniesienie kodu na energooszczędne urządzenia brzegowe z rodziny ARM.

\subsection{Komputer jednopłytkowy Raspberry Pi 5}
W segmencie urządzeń brzegowych ogólnego przeznaczenia badaniom poddano Raspberry Pi 5. W nowoczesnych topologiach Przemysłu 4.0 tego typu minikomputery pełnią zazwyczaj funkcję bramek brzegowych (IoT Gateways), agregując i filtrując dane bezpośrednio z maszyn. Wybór tej platformy podyktowany był chęcią sprawdzenia, jak "surowa", wysokotaktowana architektura ARM Cortex-A76 bez wbudowanych, dedykowanych akceleratorów sztucznej inteligencji radzi sobie ze skomplikowanymi, wielowątkowymi obciążeniami matematycznymi w języku Python. Platforma ta posłużyła do oceny balansu między przystępnością cenową i uniwersalnością sprzętu, a jego realną przepustowością w zadaniach klasycznej analizy danych \cite{rpi5specs}.

\subsection{NVIDIA Jetson Orin Nano (Edge AI)}
Trzecim środowiskiem testowym był moduł NVIDIA Jetson Orin Nano. W przeciwieństwie do Raspberry Pi, jest to urządzenie wysoce wyspecjalizowane, zaprojektowane od podstaw dla systemów autonomicznych i przemysłowych, gdzie priorytetem jest niezawodność procesora oraz dostęp do zunifikowanej pamięci operacyjnej współdzielonej między CPU i GPU. Jetson został wybrany do zestawienia jako reprezentant układów ukierunkowanych na sztuczną inteligencję. Ewaluacja na tym urządzeniu miała charakter krytyczny. Jej celem było zbadanie, jak asymetryczna architektura sprzętowa, oparta na wolniej taktowanym, ale bezpiecznym procesorze z rodziny Safety Ready oraz potężnym GPU, reaguje na obciążenie klasycznymi algorytmami CPU-bound jak drzewa decyzyjne z pakietu scikit-learn czy river, których kod źródłowy uniemożliwia wydelegowanie operacji na rdzenie CUDA \cite{jetsonspecs}.

\subsection{Mikrokontroler ESP32 – Warstwa czujnikowa (Deep Edge)}

Ostatnim, a zarazem najbardziej ograniczonym sprzętowo środowiskiem był mikrokontroler ESP32. W architekturach Internetu Rzeczy (IoT) układy tego typu reprezentują tzw. głęboki brzeg sieci (Deep Edge), pełniąc zazwyczaj rolę węzłów czujnikowych. W przeciwieństwie do wcześniej opisanych platform minikomputerowych, ESP32 nie dysponuje pełnoprawnym systemem operacyjnym ani gigabajtami pamięci współdzielonej, operując zasobami RAM rzędu zaledwie kilkuset kilobajtów oraz niskotaktowanym mikrokontrolerem.
Włączenie tego urządzenia do zestawienia miało na celu weryfikację limitów optymalizacji oprogramowania oraz możliwości wdrożenia paradygmatu TinyML (uczenia maszynowego na mikrokontrolerach) \cite{capogrosso2024tinyml}. ESP32 posłużył do przetestowania dedykowanego mechanizmu przeniesienia modeli, wyuczonych uprzednio w środowisku PC w postaci surowych struktur danych wyeksportowanych do pamięci stałej (Flash) mikrokontrolera.


\section{Analiza i porównanie platform sprzętowych oraz architektur procesorów}
Kluczowym elementem oceny wydajności algorytmów na różnych platformach (Edge vs. PC) jest dogłębna analiza ich specyfikacji sprzętowej. W niniejszym projekcie zastosowano asymetryczny podział zadań, co wymusiło ewaluację urządzeń w dwóch odmiennych kontekstach operacyjnych. Do fazy trenowania i profilowania modeli (wymagającej znacznych zasobów RAM i mocy obliczeniowej) wykorzystano trzy urządzenia: klasyczny komputer osobisty pełniący rolę punktu odniesienia, wszechstronny komputer jednopłytkowy Raspberry Pi 5 oraz wysoce wyspecjalizowany moduł sztucznej inteligencji NVIDIA Jetson Orin Nano. Z kolei do fazy czystego wnioskowania (ang. \textit{inference}) włączono skrajnie ograniczony zasobowo mikrokontroler z rodziny ESP32. Zestawienie podstawowych parametrów badanych procesorów przedstawiono w Tabeli \ref{tab:cpu}. \cite{rpi5specs}\cite{jetsonspecs}\cite{amdspecs}\cite{esp32}

\begin{table}[htpb]
    \centering
    \caption{Zestawienie parametrów technicznych wykorzystanych platform obliczeniowych}
    \label{tab:cpu}
    \renewcommand{\arraystretch}{1.3} % Zwiększa odstępy między wierszami dla lepszej czytelności
    % Zmodyfikowane szerokości kolumn, aby pomieścić 5 elementów na stronie formatu A4
    \begin{tabular}{|p{2.5cm}|p{2.8cm}|p{3.5cm}|p{3.0cm}|p{2.5cm}|}
        \hline
        \textbf{Cecha} & \textbf{Raspberry Pi 5} & \textbf{NVIDIA Jetson Orin Nano} & \textbf{PC (Lenovo LOQ)} & \textbf{ESP32-WROOM} \\ 
        \hline
        \textbf{Model procesora} & Broadcom BCM2712 & Niestandardowy układ SoC NVIDIA & AMD Ryzen 7 7435HS & Espressif ESP32 \\ 
        \hline
        \textbf{Architektura} & Arm Cortex-A76 (64-bit) & Arm Cortex-A78AE (64-bit) & x86\_64 (Zen 3+) & Xtensa LX6 (32-bit) \\ 
        \hline
        \textbf{Rdzenie / Wątki} & 4 / 4 & 6 / 6 & 8 / 16 & 2 / 2 \\ 
        \hline
        \textbf{Taktowanie} & 2,4 GHz & 1,5 GHz & 3,1 GHz (Boost do 4,5 GHz) & 240 MHz \\ 
        \hline
        \textbf{Pamięć RAM} & 4GB LPDDR4X & 8GB LPDDR5 (zintegrowana) & Obsługa DDR5 (niezależna) & 520 KB SRAM \\ 
        \hline
        \textbf{Pobór mocy (TDP)} & Bardzo niski (ok. 5-12W) & Niski (konfigurowalny 7W - 15W) & Wysoki (45W+) & Skrajnie niski ($<$ 1W) \\ 
        \hline
    \end{tabular}
\end{table}

Porównywane platformy sprzętowe reprezentują cztery zupełnie odmienne filozofie projektowania układów, co bezpośrednio determinuje ich potencjał obliczeniowy oraz obszary zastosowań. Zasadnicze różnice obejmują surową wydajność obliczeniową niezbędną do uczenia maszynowego oraz mechanizmy zapewniające bezpieczeństwo i niezawodność.

\subsection{Wydajność obliczeniowa i architektura rdzeni}
Pod względem czystej specyfikacji technicznej najwyższym potencjałem dysponuje procesor AMD Ryzen 7 7435HS. Jest to rozbudowany układ oparty na architekturze x86\_64 (mikroarchitektura Zen 3+ Rembrandt R), wyposażony w 8 rdzeni fizycznych z obsługą 16 wątków. Charakteryzuje się on najszybszym zegarem, taktowanie bazowe wynosi 3,1 GHz, a w trybie Boost osiąga 4,5 GHz. Architekturę tę wspiera niezwykle pojemna hierarchia pamięci podręcznej (16 MB współdzielonej pamięci L3 oraz 4 MB pamięci L2), co predysponuje ten układ do błyskawicznego przetwarzania dużych wolumenów danych.

W segmencie urządzeń wbudowanych (ARM) różnice są równie istotne. Komputer jednopłytkowy Raspberry Pi 5 opiera się na układzie Broadcom BCM2712 z 4-rdzeniowym procesorem Arm Cortex-A76. Dzięki taktowaniu na poziomie 2,4 GHz oraz 2 MB współdzielonej pamięci L3, architektura ta oferuje 2- do 3-krotny wzrost wydajności względem poprzedniej generacji platformy, stanowiąc bardzo wydajną jednostkę ogólnego przeznaczenia.

Z kolei NVIDIA Jetson Orin Nano (w wariancie 8 GB RAM) wykorzystuje 6-rdzeniowy procesor w nowszej architekturze Arm Cortex-A78AE. Choć jego taktowanie wynosi zaledwie 1,5 GHz, sama struktura rdzenia A78AE oferuje o 30 wyższą wydajność przetwarzania instrukcji na cykl w stosunku do starszych generacji układów ARM. Projektanci firmy NVIDIA postawili w tym przypadku na ekstremalną efektywność energetyczną oraz zoptymalizowanie potoku pod kątem równoległych operacji sztucznej inteligencji, świadomie rezygnując z wysokich częstotliwości taktowania rdzeni CPU.

\subsection{Niezawodność i bezpieczeństwo sprzętowe}
Analizowane układy różnią się drastycznie pod kątem mechanizmów bezpieczeństwa, co wynika z ich rynkowego przeznaczenia:
\begin{enumerate}[label=\alph*), leftmargin=1.25cm] 
	\item \textbf{Cortex-A78AE (NVIDIA Jetson)} -- Jest to architektura zbudowana wokół koncepcji "Safety Ready", dedykowana do systemów autonomicznych, robotyki i przemysłu ciężkiego. Jej najważniejszą cechą jest sprzętowa obsługa funkcji Split-Lock. Pozwala ona rdzeniom na pracę w trybie wysokiej niez. W tym trybie dwa rdzenie wykonują synchronicznie tę samą operację, wzajemnie weryfikując swoje wyniki, co pozwala na natychmiastowe wykrycie i wyeliminowanie błędów obliczeniowych (np. wywołanych promieniowaniem kosmicznym)
	\item \textbf{Zen 3+ (AMD Ryzen)} -- Architektura ta skupia się przede wszystkim na stabilności i inteligentnym zarządzaniu energią w ekosystemach serwerowych i konsumenckich, wprowadzając ponad 50 nowych stanów zasilania względem poprzedniej generacji.
	\item \textbf{Cortex-A76 (Raspberry Pi 5)} -- Oferuje sprzętowe rozszerzenia kryptograficzne, które znacząco przyspieszają szyfrowanie strumieni danych (np. SSL/TLS). Stanowi to doskonałe zabezpieczenie dla standardowych aplikacji sieciowych, jednak sam procesor nie posiada certyfikacji ani mechanizmów nadmiarowych (jak Jetson) pozwalających na użycie go w krytycznych systemach podtrzymujących funkcje życiowe lub sterujących autonomicznymi pojazdami.
\end{enumerate}

\subsection{Architektura mikrokontrolera ESP32 w roli wyłącznie jednostki wnioskującej}
Całkowicie odrębną klasę sprzętu reprezentuje układ SoC Espressif ESP32, stanowiący w projekcie platformę docelową dla wyuczonych algorytmów (urządzenie wykonawcze). Został on zaprojektowany w oparciu o 32-bitową mikroarchitekturę Tensilica Xtensa Dual-Core 32-bit LX6. Brak wbudowanej jednostki zarządzania pamięcią (MMU) typowej dla układów ARM/x86 oraz zaledwie 520 KB wbudowanej pamięci SRAM uniemożliwia uruchomienie pełnoprawnego systemu operacyjnego Linux i środowiska Python, dyskwalifikując układ z zadań treningowych.

Jednakże, dzięki dwóm rdzeniom taktowanym zegarem do 240 MHz, sprzętowej jednostce zmiennoprzecinkowej (FPU) realizującej instrukcje w pojedynczej precyzji oraz mechanizmom bezpośredniego mapowania pamięci stałej Flash do przestrzeni adresowej procesora (Flash cache), architektura Xtensa jest w stanie błyskawicznie wykonywać operacje mnożenia macierzy oraz iterować przez struktury drzew decyzyjnych w języku C/C++. Taka budowa czyni ESP32 optymalną architekturą brzegową do oceny i predykcji po stronie czujnika (Deep Edge) \cite{esp32}.

\section{Podstawy teoretyczne wykorzystanych algorytmów detekcji anomalii}
W ramach niniejszego projektu do detekcji anomalii wykorzystano trzy algorytmy reprezentujące odmienne paradygmaty uczenia maszynowego: sieci neuronowe oparte na rekonstrukcji sygnału, klasyczne metody zespołowe oparte na drzewach decyzyjnych oraz algorytmy strumieniowe.
\subsection{Autoencodery}
Autoenkoder to specyficzny rodzaj sztucznej sieci neuronowej, którego celem jest rekonstrukcja danych wejściowych na warstwie wyjściowej. Sieć ta składa się z dwóch głównych części: enkodera, który kompresuje dane wielowymiarowe do ukrytej reprezentacji o niższym wymiarze, tzw. wąskiego gardła, oraz dekodera, który stara się odtworzyć pierwotny sygnał na podstawie tej skompresowanej reprezentacji. W kontekście detekcji anomalii, model ten jest trenowany wyłącznie na danych opisujących normalne działanie systemu. W rezultacie sieć uczy się ignorować szum i poprawnie rekonstruować powtarzalne wzorce. Gdy na wejściu sieci pojawi się anomalia, której model wcześniej nie widział, enkoder nie potrafi poprawnie zmapować go w swojej warstwie ukrytej, co skutkuje wysokim błędem rekonstrukcji, najczęściej mierzonym jako błąd średniokwadratowy MSE (\ref{eq:mse}). Obserwacje przekraczające wyznaczony próg błędu są klasyfikowane jako anomalie.

\begin{equation}
    MSE = \frac{1}{n} \sum_{i=1}^{n} (y_i - \hat{y}_i)^2
    \label{eq:mse}
\end{equation}

Jak wykazali Chandola i in. w klasycznym przeglądzie literatury dotyczącym detekcji anomalii \cite{anomalysurvey}, modele oparte na rekonstrukcji są niezwykle skuteczne w przypadku danych wielowymiarowych. Znajdują one szerokie zastosowanie w analizie sygnałów z wielu zintegrowanych sensorów przemysłowych, gdzie kluczowe jest wychwycenie skomplikowanych, nieliniowych zależności pomiędzy poszczególnymi czujnikami, które są trudne do zamodelowania tradycyjnymi metodami statystycznymi.

\subsection{Isolation forest}
Algorytm Lasu Izolacyjnego różni się fundamentalnie od większości metod detekcji, które starają się najpierw zdefiniować profil normalności, a następnie szukają odchyleń od tej normy. Metoda ta skupia się na bezpośrednim odizolowaniu anomalii. Algorytm buduje zespół losowych drzew binarnych. W każdym kroku losowo wybierana jest cecha, np. odczyt z konkretnego czujnika, a następnie losowo dobierany jest punkt podziału pomiędzy minimalną a maksymalną wartością tej cechy. Z racji tego, że anomalie są z definicji nieliczne i odmienne, wymagają one znacznie mniejszej liczby podziałów, aby zostać odizolowanymi od reszty zbioru danych w pojedynczym liściu drzewa. W związku z tym, ścieżka od korzenia do liścia dla anomalii jest znacznie krótsza niż dla normalnych próbek. Ocena anomalii dla każdego punktu danych jest wyliczana na podstawie średniej długości ścieżki we wszystkich wygenerowanych drzewach.

Zaproponowany przez autorów Liu, Tinga i Zhou \cite{isoforest} algorytm charakteryzuje się liniową złożonością czasową oraz niskim zapotrzebowaniem na pamięć operacyjną. Ponieważ nie opiera się na obliczaniu odległości czy gęstości, świetnie sprawdza się w wielowymiarowych zbiorach danych. Jest szeroko stosowany w cyberbezpieczeństwie oraz diagnostyce maszyn przy przetwarzaniu wsadowym, będąc wysoce odpornym na szum w danych treningowych.

\subsection{Half-Space Trees}
Algorytm Half-Space Trees to nowsze podejście oparte na strukturach drzewiastych, jednak zaprojektowane od podstaw do pracy w środowisku strumieniowym. W przeciwieństwie do standardowych algorytmów wsadowych, które wymagają wczytania całego zbioru danych do pamięci, HS-Trees uczy się inkrementalnie próbka po próbce. Algorytm tworzy zbiór drzew losowych, gdzie każda przestrzeń cech jest dzielona na równe półprzestrzenie, stąd nazwa Half-Space. Przechodzące przez system dane aktualizują liczniki w węzłach drzew. Model ten często wykorzystuje mechanizm okna czasowego, dzięki czemu starsze dane są zapominane, a algorytm adaptuje się do zmieniającego się w czasie profilu normalnego zachowania maszyny. Próbka, która trafia do węzła o bardzo niskim liczniku wystąpień w danym oknie czasowym, otrzymuje wysoką punktację anomalii.

Jak udowodniono w pracy twórców algorytmu, HS-Trees to rozwiązanie dedykowane dla środowisk, w których dane napływają w sposób ciągły z dużą prędkością, a zasoby obliczeniowe i pamięciowe są ograniczone \cite{fastanomalydetection}. ak pokazują najnowsze badania, takie lekkie algorytmy są kluczowe dla skutecznej, pokładowej detekcji anomalii bezpośrednio w sieciach sensorowych \cite{reis2025edge}. Ze względu na swój stały czas aktualizacji i minimalny odcisk w pamięci RAM, algorytm ten jest idealnym kandydatem do zastosowań w ekosystemach Internetu Rzeczy oraz platformach brzegowych, gdzie urządzenia takie jak Raspberry Pi muszą dokonywać predykcji w czasie rzeczywistym na bezpośrednich odczytach z maszyn.

\section{Przygotowanie środowiska uruchomieniowego Python na platformie Raspberry Pi}
Podczas konfiguracji platformy Raspberry Pi napotkano na trudności związane z instalacją bibliotek projektowych za pomocą globalnego menedżera pakietów pip. Próby te kończyły się krytycznym błędem informującym o "środowisku zarządzanym zewnętrznie", co jest bezpośrednim wynikiem rygorystycznego wdrożenia standardu PEP 668 w najnowszych systemach operacyjnych opartych na architekturze Debian. Standard ten wprowadzono w celu eliminacji konfliktów zależności pomiędzy systemowym instalatorem apt, a menedżerem pakietów języka Python. Ponieważ współczesne dystrybucje Linuksa opierają stabilność swoich wewnętrznych usług na konkretnych wersjach bibliotek Pythona, modyfikacja globalnej przestrzeni nazw przez użytkownika mogłaby doprowadzić do awarii całego systemu. W odpowiedzi na wykrycie znacznika EXTERNALLY-MANAGED w systemowym katalogu, narzędzie pip celowo blokuje instalację, wymuszając tym samym stosowanie bezpieczniejszych praktyk inżynierii oprogramowania.

Aby zainstalować zależności wymagane przez algorytmy uczenia maszynowego bez ryzyka naruszenia stabilności powłoki systemu, konieczne było utworzenie odizolowanego środowiska projektowego. W tym celu wykorzystano wbudowany w język Python moduł venv. W pierwszej kolejności, w docelowym katalogu projektu, wygenerowano wirtualne środowisko uruchomieniowe przy użyciu polecenia \ref{lst:venv}.

\begin{lstlisting}[language=bash, numbers=none, caption={Inicjalizacja środowiska venv}, label={lst:venv}]
python -m venv .venv
\end{lstlisting}

Komenda ta utworzyła wydzieloną strukturę katalogów, zawierającą niezależną kopię interpretera języka Python oraz własną, pustą przestrzeń na pakiety. Następnie, aby skierować system do korzystania z tej wyizolowanej instancji, środowisko zostało aktywowane komendą \ref{lst:source}.
\begin{lstlisting}[language=bash, numbers=none, caption={Włączenie środowiska venv}, label={lst:source}]
source .venv/bin/activate
\end{lstlisting}

Przejście w tryb środowiska wirtualnego, sygnalizowane pojawieniem się prefiksu .venv w wierszu poleceń, odblokowało pełny, bezpieczny dostęp do menedżera pakietów. Dopiero w tak odizolowanym kontenerze możliwe było pobranie i zainstalowanie wszystkich niezbędnych bibliotek za pomocą polecenia pip install. Zastosowana procedura nie tylko pozwoliła na sprawne ominięcie blokad nałożonych przez system operacyjny, ale również zagwarantowała, że specyficzne wersje pakietów użytych do trenowania modeli detekcji anomalii pozostaną w pełni hermetyczne i nie wpłyną negatywnie na działanie innych usług działających na urządzeniu Raspberry Pi.

\section{Przygotowanie środowiska na platformie Jetson Orin Nano}
Początkowy etap prac projektowych obejmował przygotowanie odpowiedniego środowiska operacyjnego na module NVIDIA Jetson Orin Nano, co wiązało się z nieoczekiwanymi trudnościami sprzętowymi. Pierwsza próba polegała na wykorzystaniu posiadanego dysku z zainstalowanym systemem Kali Linux. Próba ta zakończyła się natychmiastowym niepowodzeniem ze względu na fundamentalne różnice w architekturze procesorów. Powszechnie stosowane systemy operacyjne z komputerów osobistych są kompilowane pod architekturę x86\_64, która jest całkowicie niekompatybilna z układami wbudowanymi firmy NVIDIA, opartymi na architekturze ARM64. Moduł Jetson nie był w stanie rozpoznać ani uruchomić instrukcji z takiego nośnika.

Wobec powyższego, dysk został sformatowany w celu przeprowadzenia czystej instalacji dedykowanego systemu. Podjęto próbę klasycznego zbootowania instalatora z nośnika USB bezpośrednio na przygotowany dysk docelowy. To podejście również zakończyło się fiaskiem, uwydatniając kolejną cechę charakterystyczną platform brzegowych. Urządzenia z rodziny Jetson nie posiadają klasycznego układu BIOS/UEFI znanego z komputerów PC, który pozwalałby na swobodne bootowanie z zewnętrznych urządzeń USB. Ich proces rozruchu opiera się na specyficznym bootloaderze, który do inicjalizacji systemu na zewnętrznych dyskach NVMe/SSD wymaga przeprowadzenia skomplikowanego procesu flashingu za pomocą zewnętrznego komputera-hosta i narzędzia NVIDIA SDK Manager.

Ostatecznie, aby uzyskać stabilne i działające środowisko w sposób optymalny czasowo, konieczna była zmiana nośnika pamięci. Zgodnie z natywnym wsparciem rozruchowym modułu, zakupiono kartę pamięci microSD, na którą wgrano dedykowany, oficjalny obraz systemu operacyjnego Linux for Tegra JetPack OS. Wykorzystanie wbudowanego interfejsu kart pamięci pozwoliło na natychmiastowe i bezproblemowe zbootowanie systemu. Dopiero ten krok umożliwił przejście do właściwych prac związanych z konfiguracją środowiska programistycznego i testowaniem algorytmów uczenia maszynowego.


\section{Problem akceleracji GPU na platformie NVIDIA Jetson}
Podczas prób wdrożenia i przyspieszenia algorytmów detekcji anomalii na platformie NVIDIA Jetson Orin Nano, napotkano barierę programową. Pomimo fizycznej obecności wydajnego układu graficznego, żadna z wykorzystywanych bibliotek nie była w stanie wydelegować do niego obliczeń. Analiza tego zjawiska wykazała, że problem wynika bezpośrednio z braku kompatybilności ekosystemu klasycznego uczenia maszynowego z architekturą CUDA w środowisku ARM64. Ograniczenia te można podzielić na cztery główne kategorie.
\subsection{Brak natywnego wsparcia GPU w standardowych bibliotekach w języku Python}
Podstawowe biblioteki użyte w projekcie, takie jak scikit-learn, odpowiedzialna za modele Isolation Forest oraz MLPRegressor oraz river, implementująca algorytm Half-Space Trees, zostały zaprojektowane od podstaw z myślą o wykonywaniu obliczeń wyłącznie na głównym procesorze. Ich kod źródłowy nie posiada zaimplementowanych tzw. backendów dla technologii NVIDIA CUDA. Oznacza to, że niezależnie od konfiguracji sprzętowej urządzenia, biblioteki te fizycznie nie potrafią skompilować swojego grafu obliczeniowego do instrukcji zrozumiałych dla rdzeni graficznych i zawsze wymuszają wykonanie kodu na procesorze głównym.
\subsection{Ograniczenia pakietu NVIDIA RAPIDS (cuML) na architekturze ARM64}
Rozwiązaniem problemu braku wsparcia GPU dl w środowisku PC jest zazwyczaj użycie dedykowanej biblioteki NVIDIA RAPIDS (cuML). Niestety, platforma Jetson opiera się na architekturze ARM64, podczas gdy pakiety RAPIDS są rozwijane, kompilowane i optymalizowane głównie pod kątem dużych serwerów opartych na architekturze x86\_64. Próba instalacji tych bibliotek za pomocą standardowych menedżerów pakietów, np. pip na Jetsonie kończy się niepowodzeniem z powodu braku gotowych, prekompilowanych plików binarnych. Wymusza to niezwykle skomplikowaną kompilację ze źródeł, która na urządzeniach brzegowych często skutkuje konfliktami wersji kompilatora lub niekompatybilnością z wersją oprogramowania JetPack.
\subsection{Konieczność zmiany frameworku dla sieci neuronowych}
W przypadku algorytmu Autoenkodera, realizowanego za pomocą klasy MLPRegressor z biblioteki scikit-learn, ograniczenie jest identyczne, biblioteka ta nie potrafi korzystać z akceleracji sprzętowej. Aby wykorzystać potencjał modułu Jetson Orin Nano do sieci neuronowych, konieczne byłoby całkowite przepisanie architektury modelu z użyciem frameworków takich jak PyTorch lub TensorFlow. Tylko te potężne biblioteki Deep Learningowe posiadają pełne, natywne wsparcie dla ekosystemu JetPack oraz potrafią bezpośrednio mapować operacje na rdzenie CUDA i Tensor w architekturze ARM.
\subsection{Próba wdrożenia oficjalnego środowiska kontenerowego (NVIDIA L4T ML)}
W ramach pogłębionej diagnostyki i wykluczenia ewentualnych błędów konfiguracyjnych systemu operacyjnego, podjęto próbę uruchomienia skryptów wewnątrz oficjalnego kontenera Docker, udostępnianego przez firmę NVIDIA dla platform brzegowych, z rodziny Linux for Tegra Machine Learning. Chociaż to wysoce zoptymalizowane środowisko dostarcza preinstalowane sterowniki CUDA, cuDNN oraz TensorRT, nie przyniosło ono oczekiwanego rozwiązania w postaci akceleracji GPU dla badanych algorytmów. Analiza zawartości kontenera wykazała, że oficjalne obrazy NVIDIA są przygotowywane niemal wyłącznie pod kątem wsparcia głębokich sieci neuronowych. Wchodzące w skład kontenera biblioteki do klasycznej analizy danych, takie jak scikit-learn czy pakiety strumieniowe, to wciąż ich standardowe, procesorowe dystrybucje. Środowisko kontenerowe, nawet z poprawnie zmapowanym akceleratorem graficznym, nie potrafi automatycznie przetłumaczyć ani wymusić wykonania instrukcji sprzętowych dla algorytmów, które architektonicznie nie posiadają zaprogramowanego backendu CUDA. Eksperyment ten ostatecznie potwierdził, że brak akceleracji wynika bezpośrednio z braku implementacji algorytmów na GPU w ekosystemie ARM64, a nie z błędnej konfiguracji samego urządzenia Jetson.

Zjawisko braku możliwości wykonania algorytmów na GPU udowadnia, że przeniesienie projektu na platformę brzegową wymaga ścisłego doboru technologii oprogramowania. Klasyczne uczenie maszynowe w języku Python pozostaje w dużej mierze domeną procesorów. Próba wymuszenia akceleracji sprzętowej na platformie ARM dla algorytmów niebędących głębokimi sieciami neuronowymi napotyka na nieprzezwyciężone w standardowy sposób bariery w samym kodzie źródłowym dostępnych bibliotek.

\section{Wdrożenie algorytmów na mikrokontrolerze ESP32 – asymetryczny podział zadań}
Próba bezpośredniego uruchomienia zaawansowanych skryptów analitycznych na mikrokontrolerze ESP32 napotyka na bariery sprzętowe i programowe. Układ ten dysponuje zaledwie 520 KB wbudowanej pamięci operacyjnej SRAM oraz architekturą pozbawioną zaawansowanej jednostki zarządzania pamięcią MMU. Ograniczenia te uniemożliwiły w przypadku niniejszego projektu instalację pełnoprawnego systemu operacyjnego, a w konsekwencji – uruchomienie standardowego interpretera języka Python wraz z bibliotekami matematycznymi, takimi jak \textit{scikit-learn} czy \textit{river}. Wykorzystanie alternatywnych, odchudzonych środowisk pokroju MicroPython również nie było rozwiązaniem problemu. Nie wspierają one skomplikowanych frameworków uczenia maszynowego, a minimalna ilość pamięci RAM nie pozwala na wczytanie obszernych, historycznych zbiorów danych treningowych. Z tego względu przeprowadzenie pełnego cyklu uczenia modeli bezpośrednio na warstwie czujnikowej było niemożliwe do zrealizowania w ramach projektu przez problemy architektonicznie.

Aby zrealizować cel projektu i wdrożyć sztuczną inteligencję na najniższym poziomie sprzętowym, konieczne było zastosowanie paradygmatu TinyML, opierającego się na ścisłym, asymetrycznym podziale zadań. Kosztowną obliczeniowo fazę profilowania i uczenia algorytmów oddelegowano do środowiska referencyjnego PC w całości. Następnie, aby przenieść uzyskaną wiedzę na mikrokontroler, konieczne było całkowite porzucenie środowiska Python na rzecz kompilowanego języka C/C++. Wyuczone w Pythonie modele zostały wyekstrahowane z dynamicznych struktur obiektowych do postaci płaskich, statycznych tablic danych jako macierze wag dla sieci neuronowych oraz zbiory struktur węzłów dla drzew decyzyjnych. Zapisano je w plikach nagłówkowych (format \texttt{.h}) i wkompilowano bezpośrednio do pamięci stałej Flash układu ESP32. W rezultacie mikrokontroler nie buduje modeli matematycznych od podstaw. ESP32 w ramach projektu realizuje wyłącznie fazę wnioskowania (ang. \textit{inference}). Mechanizm ten, polegający na przepuszczaniu nowych sygnałów z czujników przez z góry zdefiniowane reguły i tablice wag za pomocą podstawowych operacji arytmetycznych, pozwala na szybkie i deterministyczne wykrywanie anomalii przy niewielkim zużyciu zasobów. Pozwoliło to na całkowite ominięcie problem braku środowiska Python na urządzeniu docelowym.

\section{Charakterystyka wykorzystanego zbioru danych}
Do celów badawczych, treningu oraz ewaluacji zaimplementowanych algorytmów wykorzystano ogólnodostępny zbiór danych Pump Sensor Data udostępniony na platformie Kaggle \cite{kaggledata}. Reprezentuje on rzeczywiste, historyczne odczyty parametrów pracy przemysłowej pompy wodnej zlokalizowanej w niewielkiej miejscowości. Zbiór ten stanowi klasyczny przykład danych typu Internet of Things pochodzących z systemów SCADA, co czyni go idealnym fundamentem do symulacji obciążeń na platformach brzegowych. Struktura zbioru podizelona jest na trzy główne grupy informacyjne.
\subsection{Przestrzeń czasowa (Timestamp Data)}
Zbiór określa dokładny moment wykonania każdego pomiaru. Odczyty ze wszystkich czujników były rejestrowane synchronicznie w stałych, jednominutowych interwałach. Obejmują one okres 153 dni pracy pompy, co łącznie przekłada się na potężną bazę liczącą 220 320 unikalnych rekordów. Taka ziarnistość pozwala na traktowanie zgromadzonych informacji jako wielowymiarowych szeregów czasowych. Dzięki precyzyjnym znacznikom czasowym możliwe jest sekwencyjne odtwarzanie danych w pętli programu, symulując strumieniowanie w czasie rzeczywistym dla algorytmu Half-Space Trees oraz prowadzenie analizy trendów środowiskowych poprzedzających wystąpienie awarii.
\subsection{Przestrzeń atrybutów fizycznych (Sensor Data)}
Główną część zbioru stanowią surowe odczyty fizyczne z urządzenia. Baza zawiera 52 niezależne serie danych pochodzące z oddzielnych jednostek sensorowych maszyny, oznaczonych odpowiednio od sensor\_00 do sensor\_51. Ponieważ są to wartości surowe, prosto z systemów pomiarowych, charakteryzują się one obecnością naturalnego szumu pomiarowego, wartości odstających oraz brakami w danych. Skrajnym przypadkiem jest zmienna sensor\_15, która składa się całkowicie z wartości pustych. Taka struktura bazy wymusiła w początkowym etapie przeprowadzania badań zastosowanie mechanizmów wstępnego przetwarzania danych, w tym imputacji braków metodami propagacji wartości  oraz standaryzacji zakresów za pomocą algorytmu MinMaxScaler.
\subsection{Zmienna celu}
Kluczowym atrybutem zbioru, pełniącym funkcję referencyjną, tzw. ground truth, jest etykieta aktualnego stanu maszyny. Choć ewaluowane algorytmy należą do rodziny modeli nienadzorowanych lub częściowo nadzorowanych, etykieta ta była niezbędna do końcowej oceny ich skuteczności (wyliczenia macierzy pomyłek i metryki F1-Score). Występuje ona w trzech stanach:
\begin{enumerate}[label=\alph*), leftmargin=1.25cm] 
	\item \textbf{NORMAL (Praca normalna)} -- Oznacza poprawne funkcjonowanie pompy. W ramach metodyki projektowej, podzbiór danych oznaczony tym statusem został odizolowany i posłużył jako jedyne źródło wiedzy dla modeli w fazie treningu (uczenie się zdrowego profilu maszyny).
	\item \textbf{BROKEN (Awaria)} -- Oznacza krytyczny moment wystąpienia usterki systemu. W analizowanym rocznym okresie badawczym odnotowano dokładnie 7 takich zdarzeń. Stanowią one główne anomalie, które implementowane modele miały za zadanie wykryć wyłącznie na podstawie odchyleń od normy w odczytach z sensorów.
	\item \textbf{RECOVERING (Okres regeneracji)} -- Wskazuje na okres przejściowy następujący po awarii, w którym system pompowy był naprawiany, stabilizowany i przywracany do pełnej sprawności operacyjnej.
\end{enumerate}
Zastosowanie tak zdefiniowanego, surowego i wysoce niezbalansowanego zbioru danych pozwoliło na przeprowadzenie realistycznych testów zaimplementowanych architektur detekcji anomalii, wiernie oddając warunki panujące w nowoczesnym przemyśle.

\section{Metodyka przygotowania danych oraz profilowania wydajnościowego}
\subsection{Akwizycja i preprocessing danych telemetrycznych}
Eksperymenty przeprowadzono z wykorzystaniem zbioru Pump sensor data, zawierającego wielowymiarowe szeregi czasowe z 52 czujników przemysłowych, rejestrowane z częstotliwością jednej minuty. W pierwszej fazie preprocessingu dokonano selekcji cech, odrzucając metadane niebędące odczytami fizycznymi (znacznik czasu, indeks, tekstowa etykieta statusu).

Ze względu na niekompletność strumienia danych (szczególnie zauważalną w przypadku sensora nr 15), zastosowano kaskadową strategię imputacji braków. W pierwszej kolejności wykorzystano metodę propagacji ostatniej znanej wartości w przód (ang. forward fill), następnie w tył (backward fill), a wszelkie rezydualne braki zastąpiono zerami.

Proces binaryzacji etykiet przekształcił problem w zadanie klasyfikacji dwuklasowej:
\begin{itemize}[label={--},labelsep=0.4cm,leftmargin=0.6cm]
\item Stan NORMAL (poprawna praca układu) zamapowano jako 0,
\item Stany awaryjne BROKEN oraz przejściowe RECOVERING zagregowano do wspólnej klasy anomalii oznaczonej jako 1.
\end{itemize}

Podziału na zbiór treningowy (70\%) i testowy (30\%) dokonano z rygorystycznym zachowaniem chronologii szeregu czasowego. Wyłączenie tasowania próbek było krytyczne, aby zapobiec wyciekowi danych z przyszłości do zbioru uczącego. Indeksy podziału utrwalono w plikach train.txt i test.txt, gwarantując pełną reprodukowalność eksperymentów dla różnych algorytmów.

Następnie surowe odczyty poddano normalizacji za pomocą transformacji MinMaxScaler, skalując przestrzeń cech do przedziału $[0, 1]$. Standaryzacja ta jest niezbędnym warunkiem stabilnej konwergencji sieci neuronowych (Autoenkoderów). W ostatnim kroku, bazując na paradygmacie uczenia nienadzorowanego, na zbiór treningowy nałożono maskę logiczną, filtrując wyłącznie odczyty zdrowej maszyny. Dzięki temu algorytmy modelowały wyłącznie profil "normy", kwalifikując wszelkie odchylenia w fazie testowej jako anomalię.

\subsection{Profilowanie termiczne i ograniczenia platformowe}
Istotnym elementem oceny wydajności modeli na urządzeniach brzegowych było monitorowanie temperatury układów obliczeniowych. Próba ujednoliconego pomiaru termicznego za pomocą skryptu napisanego w języku Python uwypukliła jednak wyzwania związane z fragmentacją środowisk sprzętowych i architekturą systemów operacyjnych. W toku eksperymentów ustalono, że stabilne i wiarygodne profilowanie temperatury udało się zrealizować wyłącznie na jednej z czterech badanych platform. Ograniczenia wynikały z następujących różnic architektonicznych:
\begin{enumerate}[label=\alph*), leftmargin=1.25cm]
\item \textbf{ESP32}
\item \textbf{Komputer osobisty oraz Nvidia Jetson} -- Na obu tych platformach próby odczytu temperatury z poziomu skryptu zakończyły się niepowodzeniem. Wynikało to ze specyfiki systemów operacyjnych, zabezpieczeń izolujących procesy użytkownika oraz braku ustandaryzowanego interfejsu takiego jak np. uniwersalny plik wirtualny w systemach Linux. Skrypt w języku Python nie był w stanie uzyskać bezpośredniego dostępu do telemetrii sensorów termicznych na płycie głównej bez głębokiej ingerencji w uprawnienia administracyjne (np. wywoływania zewnętrznych narzędzi systemowych typu lm-sensors, co zaburzałoby uniwersalność kodu).
\item \textbf{Raspberry Pi} -- Architektura systemu operacyjnego tego mikrokomputera okazała się w pełni kompatybilna z przyjętym modelem pomiarowym. Platforma ta jako jedyna pozwoliła na stabilny i ciągły odczyt temperatury układu, co zrealizowano poprzez bezpośrednie odczytywanie wartości z systemowego pliku wirtualnego /sys/class/thermal/thermal\_zone0/temp. Zebrane w ten sposób dane umożliwiły precyzyjną ocenę nagrzewania się urządzenia pod obciążeniem algorytmów analitycznych.
\end{enumerate}

\subsection{Pomiar czasów globalnych}
W tym trybie wykorzystano standardową funkcję systemową time.time(), która rejestruje bezwzględny czas wykonywania się głównych bloków kodu. Funkcję wywoływano (oraz obliczano różnicę czasu) w następujących momentach działania aplikacji:
\begin{enumerate}[label=\alph*), leftmargin=1.25cm]
    \item \textbf{Czas przygotowania danych} -- Wyznaczenie czasu potrzebnego na załadowanie zbioru CSV do pamięci, uzupełnienie braków (tzw. kaskadowy fill) oraz skalowanie cech.
    \item \textbf{Czas treningu (t\_train)} -- Znaczniki czasu wstawiono bezpośrednio przed i po bloku uczącym dany model, co pozwoliło określić globalny czas potrzebny na wytworzenie profilu normalnego zachowania na konkretnej platformie.
    \item \textbf{Całkowity czas detekcji (t\_detect\_total)} -- Pomiar ten obejmował przepuszczenie całego zbioru testowego przez nauczony model naraz, co symuluje wydajność klasycznego przetwarzania paczkowego (ang. batch processing).
\end{enumerate}

\subsection{Analiza opóźnień (Latency) i przepustowości (Throughput)}
Kluczowym kryterium ewaluacji systemów wbudowanych jest czas reakcji na nowe zdarzenia. W celu rzetelnego zmierzenia opóźnień wnioskowania (ang. inference latency), ze zbioru testowego wyizolowano losową próbę 200 obserwacji.

Oprogramowanie przetwarzało te próbki sekwencyjnie w pętli (symulując napływający strumień danych z czujników), rejestrując czas z mikrosekundową precyzją bezpośrednio przed i po wywołaniu predykcji modelu. Zebrane czasy posłużyły do wyznaczenia kluczowych metryk brzegowych:
\begin{enumerate}[label=\alph*), leftmargin=1.25cm]
\item \textbf{Średnia latencja (lat\_mean)} – Przeciętny czas, jakiego model potrzebuje na przeanalizowanie jednego wektora cech (odpowiadającego jednej minucie pracy maszyny).
\item \textbf{Kwantyle opóźnień (lat\_p95, lat\_p99)} – Wartości określające górną granicę czasu potrzebnego na przetworzenie odpowiednio 95\% i 99\% najgorszych (najdłużej analizowanych) próbek. Metryka P99 jest w inżynierii standardem określającym stabilność i gwarancję czasu odpowiedzi (SLA) w systemach czasu rzeczywistego.
\item \textbf{Przepustowość (throughput)} – Wyrażona w liczbie przetworzonych próbek na sekundę, określa maksymalną wydajność systemu. Informuje, ile minut historycznej lub napływającej telemetrii dane urządzenie jest w stanie przeanalizować w ciągu jednej sekundy czasu rzeczywistego.
\end{enumerate}

\subsection{Profilowanie obciążenia procesora}
Do monitorowania aktywności jednostki obliczeniowej wykorzystano zewnętrzną bibliotekę psutil, skonfigurowaną do pracy w trybie różnicowym (z parametrem interval=None). Podejście to pozwala na asynchroniczne badanie średniego obciążenia procesora pomiędzy zdefiniowanymi punktami kontrolnymi w kodzie, bez sztucznego wstrzymywania i blokowania głównego wątku aplikacji. Cykl zbierania metryk podzielono na trzy strategiczne etapy:
\begin{enumerate}[label=\alph*), leftmargin=1.25cm]
\item \textbf{Inicjalizacja punktu odniesienia} -- Pierwsze wywołanie funkcji na starcie skryptu, pełniące rolę kalibracji systemu (uruchomienie wewnętrznego timera biblioteki psutil).
\item \textbf{Odczyt po przygotowaniu danych} -- Rejestracja narzutu obliczeniowego wygenerowanego przez operacje wejścia/wyjścia (I/O), proces kaskadowego uzupełniania braków oraz normalizację cech.
\item \textbf{Odczyt po sfinalizowaniu treningu} -- Ewaluacja kosztu zasobów zużytych stricte na optymalizację algorytmu (np. budowę drzew binarnych lub dopasowanie wag w sieci neuronowej).
\end{enumerate}
Końcowy wskaźnik utylizacji procesora, prezentowany w raporcie podsumowującym, wyznaczany jest jako średnia arytmetyczna z pomiarów zarejestrowanych w etapach drugim i trzecim. Przyjęta metodyka gwarantuje miarodajną i obiektywną ocenę całkowitego kosztu obliczeniowego, jakiego wymagała kompleksowa realizacja głównej logiki detekcyjnej systemu.

\subsection{Obserwacja alokacji pamięci RAM}
Pomiary zużycia pamięci operacyjnej zrealizowano przy użyciu wbudowanej biblioteki \kod{tracemalloc}, która umożliwia precyzyjne śledzenie alokacji pamięci przez procesy środowiska Python. Proces monitorowania inicjowany jest poprzez wywołanie funkcji \texttt{tracemalloc.start()} bezpośrednio przed rozpoczęciem fazy treningowej modelu. Od tego momentu system rejestruje każdy zaalokowany blok pamięci, pozwalając na dokładne uchwycenie rzeczywistego narzutu pamięciowego algorytmu w trakcie jego pracy.

Po sfinalizowaniu procesu uczenia następuje wywoływanie kolejnej metody o nazwie \kod{tracemalloc.get\_traced\_memory(}). Funkcja ta zwraca krotkę zawierającą dwie zmienne: zużycie bieżące oraz zużycie szczytowe. Do celów analitycznych ekstrahowana jest wyłącznie wartość szczytowa, ponieważ to ona determinuje maksymalne zapotrzebowanie modelu na zasoby, co stanowi parametr krytyczny przy wdrażaniu rozwiązań na urządzenia brzegowe o ograniczonej pojemności pamięci SRAM/RAM. Uzyskany wynik, domyślnie wyrażony w bajtach, poddawany jest konwersji na megabajty poprzez podzielenie przez $1024^2$. Cykl profilowania kończy się wywołaniem \kod{tracemalloc.stop()}, co skutecznie czyści bufory monitorujące i zwalnia zasoby systemowe przypisane do analizatora.

\section{Szczegóły implementacyjne i dobór hiperparametrów modeli do procesu uczenia}
Zrozumienie działania badanych algorytmów wymaga analizy ich implementacji w kodzie oraz uzasadnienia doboru kluczowych hiperparametrów. Poniżej przedstawiono szczegółowy opis konfiguracji każdego z modeli wraz z listingami najważniejszych fragmentów kodu.
\subsection{Powtarzalność wyników badań}
We wszystkich implementowanych modelach oraz funkcjach podziału zbioru danych zastosowano stałe ziarno losowości \kod{random\_state=42} (lub \kod{seed=42} dla biblioteki river). W uczeniu maszynowym inicjalizacja wag sieci neuronowych, losowanie cech w drzewach decyzyjnych czy mieszanie próbek opierają się na generatorach liczb pseudolosowych. Ustawienie stałego ziarna, w tym przypadku popularnej w środowisku programistycznym liczby 42, jest kluczową praktyką badawczą. Gwarantuje to determinizm eksperymentu, czyli pozwala na dokładne i sprawiedliwe porównanie czasów wykonania oraz metryk F1-Score na różnych platformach sprzętowych, mając pewność, że na każdym urządzeniu zainicjalizowano matematycznie identyczny model.
\subsection{Implementacja Autoenkodera}
\begin{lstlisting}[language=python, numbers=none, caption={Parametry uczenia Autoencodera}, label={lst:autoendocer}]
# Inicjalizacja modelu
model = MLPRegressor(hidden_layer_sizes=(32, 16, 32), activation='relu', solver='adam', max_iter=300, random_state=42)

# Trening (tylko na klasie NORMAL)
mask_normal = (y_train == 0).values
X_train_normal_scaled = X_train_scaled[mask_normal]
model.fit(X_train_normal_scaled, X_train_normal_scaled)

\end{lstlisting}

Wykorzystano wielowarstwowy perceptron (\kod{MLPRegressor}), ponieważ w trybie uczenia nienadzorowanego wejście sieci (zmienna X) jest podawane również jako oczekiwane wyjście (zmienna y). Architektura składająca się z warstw o rozmiarach 32, 16 i 32 (\ref{lst:autoendocer}) węzłów realizuje fundamentalne założenie autoenkodera, tworząc tzw. "wąskie gardło" (bottleneck) w środkowej warstwie. Wejściowy wektor cech jest kompresowany z n-wymiarów początkowych przez warstwę 32 neuronów do zaledwie 16 neuronów, a następnie rekonstruowany z powrotem przez warstwę 32 neuronów do oryginalnego wymiaru. Taka kompresja zapobiega zjawisku, w którym sieć uczy się zwykłej funkcji tożsamościowej, czyli przepisywania wejścia na wyjście 1:1 i zmusza model do wyekstrahowania najważniejszych, ukrytych cech sygnału.

Linijki izolujące klasę \kod{NORMAL} i trenujące model wyłącznie na niej stanowią esencję detekcji anomalii. Gdybyśmy do modelu wprowadzili cały zbiór, łącznie z anomaliami, sieć nauczyłaby się rekonstruować również błędy. Ucząc ją tylko poprawnych zachowań maszyny, zapewniamy, że błąd rekonstrukcji MSE wystrzeli w górę, gdy pojawi się nieznana anomalia.

\subsection{Implementacja Isolation Forest}

\begin{lstlisting}[language=python, numbers=none, caption={Parametry uczenia Isolation Forest}, label={lst:forest}]
model = IsolationForest(n_estimators=100, n_jobs=-1, random_state=42)
model.fit(X_train_normal_scaled)

# Odwrocenie wynikow funkcji decyzyjnej
test_scores = -model.decision_function(X_test_scaled)
\end{lstlisting}

Liczba 100  (Listing \ref{lst:forest}) określa rozmiar lasu (ilość drzew izolacyjnych). W oryginalnej pracy badawczej dotyczącej tego algorytmu autorzy udowodnili, że długości ścieżek w drzewach decyzyjnych bardzo szybko zbiegają się do ostatecznych wartości. Zwiększenie tej liczby np. do 500 czy 1000 drzew drastycznie wydłużyłoby czas obliczeń i zużycie pamięci RAM na urządzeniach brzegowych (Raspberry Pi, Jetson), nie przynosząc zauważalnej poprawy trafności klasyfikacji. Wartość 100 stanowi optymalny kompromis między wydajnością a szybkością.

Zastosowanie operacji \kod{-model.decision\_function()} jest niezwykle ważnym zabiegiem inżynierskim w kodzie. Domyślnie biblioteka \kod{scikit-learn} zwraca wyniki ujemne dla próbek anomalnych, a dodatnie dla normalnych. Odwrócenie znaku wartości ujednolica logikę programu dla wszystkich trzech algorytmów, w każdym przypadku wprowadzona zostaje zasada “im wyższy wynik punktowy, tym większa anomalia", co pozwala na zastosowanie jednej uniwersalnej funkcji wyliczania progu.

\subsection{Implementacja Half-Space Trees}
\begin{lstlisting}[language=python, numbers=none, caption={Parametry uczenia HS-Trees}, label={lst:hs}]
model = anomaly.HalfSpaceTrees(n_trees=25, height=15, window_size=70000, seed=42)

# Petla strumieniowego uczenia online
for i, xi in enumerate(X_train_normal_scaled):
    model.learn_one(dict(enumerate(xi)))

\end{lstlisting}

Algorytm strumieniowy musi reagować błyskawicznie na pojedyncze zdarzenia. Ponieważ struktura drzew HS-Trees jest aktualizowana i przeszukiwana przy każdej iteracji (dla każdej pojedynczej próbki danych), mniejsza liczba drzew, \kod{n\_trees=25}, optymalizuje opóźnienie (latencję) na słabych procesorach typu ARM. Dodatkowo parametr \kod{height=15} narzuca twardy limit na głębokość każdego drzewa (maksymalnie $2^{15}$ liści na drzewo). Zabezpiecza to moduły IoT, posiadające ograniczoną ilość pamięci operacyjnej, przed niekontrolowanym rozrostem struktur danych w przypadku długotrwałego nasłuchiwania na dane ze strumienia.

Fundamentalną cechą algorytmów strumieniowych (Online Machine Learning), odróżniającą je od klasycznych modeli wsadowych (batch processing), jest zdolność do ciągłej adaptacji i usuwania nieaktualnych danych. Biblioteka \kod{river} realizuje tę funkcjonalność za pomocą mechanizmu przesuwnego okna (sliding window). Dobór parametru \kod{window\_size=70000} w implementacji algorytmu Half-Space Trees nie był przypadkowy i wynikał bezpośrednio z charakterystyki analizowanego zbioru danych. Ponieważ całkowity wolumen użytych danych historycznych wynosił 140 000 rekordów, ustawienie szerokości okna na 70 000 oznacza, że model w każdym momencie ewaluacji przetwarza i opiera swoje decyzje na rozkładzie statystycznym stanowiącym dokładnie połowę całkowitego "cyklu życia" badanej maszyny.

Taka konfiguracja stanowi optymalny kompromis inżynierski. Z jednej strony, utrzymywanie w pamięci aż 70 000 punktów referencyjnych buduje wysoce stabilny profil "normalności". Dzięki temu system wykazuje odpowiednią bezwładność, jest odporny na chwilowe zaszumienia sygnału z sensorów i nie generuje fałszywych alarmów przy niegroźnych mikroskokach. Z drugiej strony, zastosowanie sztywnego limitu bufora zabezpiecza model przed kluczowym w systemach przemysłowych zjawiskiem dryfu koncepcji (concept drift). Układy mechaniczne ulegają naturalnemu zużyciu, a ich domyślne parametry pracy (np. drgania, temperatura) mogą z biegiem czasu powoli ewoluować. Poprzez sukcesywne usuwanie najstarszych obserwacji wykraczających poza połowę cyklu danych, algorytm HS-Trees w sposób płynny i autonomiczny adaptuje się do nowych warunków pracy, nie traktując ich od razu jako krytycznej awarii. Ponadto, narzucenie twardego limitu elementów zapamiętywanych przez struktury drzewiaste gwarantuje stałą i przewidywalną złożoność pamięciową, co chroni wbudowane platformy brzegowe przed wyczerpaniem zasobów RAM, niezależnie od tego, jak długo urządzenie pozostaje uruchomione.

Zastosowanie funkcji \kod{model.learn\_one(dict(enumerate(xi)))} uwidacznia różnicę architektoniczną biblioteki \kod{river}. Zamiana wektora \kod{NumPy} na słownik języka Python w każdej iteracji symuluje odebranie pojedynczego pakietu JSON/słownikowego bezpośrednio z maszyny przemysłowej za pośrednictwem sieci (np. protokołem MQTT), co stanowi doskonałe odwzorowanie produkcyjnego scenariusza w urządzeniach Edge.

\subsection{Obliczanie progu anomalii}
\begin{lstlisting}[language=python, numbers=none, caption={Próg anomalii}, label={lst:th}]
threshold = np.percentile(train_scores, 98.8)
pred_labels_test = (test_scores > threshold).astype(int)
\end{lstlisting}

Próg anomalii dla poszczególnych algorytmów był pozyskiwany w sposób następujący:
\begin{enumerate}[label=\alph*), leftmargin=1.25cm] 
\item \textbf{Autoencoder} --
poprzez funkcje model.predict(\kod{X\_train\_normal\_scaled}) generowana jest odpowiedź modelu jak wyglądają dane wyjściowe na próbach na których model był uczony. Następnie liczony jest średni błąd rekonstrukcji (różnica między danymi wejściowymi w wiśniowymi, to podnosi kwadratu). Bład został zapisany do zmiennej train\_scores
\item \textbf{Isolation forest} --
fukcja model.decision\_function zwraca im niższa wartość im krótsza jest ścieżka, aby logika pasowała do reszty kodu dodano na początku tej funkcji minus aby ta logikę odwrócić i aby krótsza ścieżka otrzymywała wyższa wartość wartość. Następnie wartość zapisywana jest do zmiennej train\_scores 

\item \textbf{Half-space tree} -- 
Funkcja score one wylicza wynik na podstawie tego, jak wiele normalnych próbek wylądowało wcześniej w tych samych obszarach co badana próbka, obliczone wartości  trafiają do wartość train\_score.  

\end{enumerate}
Następnie liczony jest próg anomalii thershold. Używana jest funkcja percentyla, która sortuje wszystkie liczby w tablicy od najmniejszego do największego. Próg zostaje ustawiony w takim miejscu, że 98.8\% wszystkich normalnych odczytów znalazło się poniżej niego. Pozostałe 1.2\% wyników (najwyższe błedy w zbiorze normalnym) jest traktowanej jako anomalia.

\section{Proces klasyfikacji i detekcji na zbiorze testowym}
Po ustaleniu granicy decyzyjnej (threshold) na podstawie danych treningowych, system przechodzi do docelowego etapu klasyfikacji na niewidzianym wcześniej zbiorze testowym. Ewaluacja ta przebiega symetrycznie do fazy uczenia. W pierwszej kolejności modele analizują nowe próbki i generują surowe miary odchyleń, np. błędy rekonstrukcji, długości ścieżek drzew binarnych, które są agregowane w wektorze wynikowym o nazwie \kod{test\_scores}.

Następnie implementowana jest właściwa logika detekcyjna. Każda wyestymowana wartość numeryczna z wektora testowego poddawana jest weryfikacji relacyjnej względem wcześniej wyliczonego progu anomalii. Algorytm aplikuje na zbiorze warunek boolowskim, jeśli wynik punktowy danej próbki jest ostro większy niż założony próg, operacja zwraca wartość logiczną \kod{True}. Jeśli natomiast odchylenie jest mniejsze lub równe wartości progowej, układ interpretuje to jako pracę w normie, zwracając \kod{False}.

W celu ujednolicenia formatu danych z matematycznymi standardami ewaluacji (wymaganymi do wyliczenia macierzy pomyłek oraz metryk precyzji), uzyskany wektor logiczny jest natychmiastowo poddawany rzutowaniu typów. Za pomocą wbudowanej metody konwersji (\kod{.astype(int)}), wartości logiczne zamieniane są na reprezentację całkowitoliczbową. W rezultacie powstaje finalna tablica predykcji o nazwie \kod{pred\_labels\_test}. Składa się ona wyłącznie z wartości binarnych, zera (0) oznaczającego poprawny stan maszyny (NORMAL) oraz jedynki (1) jednoznacznie wskazującej na zidentyfikowaną usterkę (ANOMALIA). Tablica ta stanowi bezpośredni materiał wejściowy do oceny skuteczności klasyfikacyjnej sprzętu.

\section{Implementacja detekcji anomalii na module ESP32}

W tym rozdziale przedstawiono wdrożenie algorytmów na docelowym mikrokontrolerze ESP32. Ze względu na ograniczenia sprzętowe, deficyt pamięci SRAM, wykluczono możliwość uruchomienia środowiska Python oraz ciężkich bibliotek analitycznych. Przyjęto założenie opracowania specyficznej reprezentacji modeli. Wiedza wyuczona w środowisku PC (wagi, obciążenia oraz parametry węzłów) została zmieniona do postaci statycznych tablic i ulokowana w pamięci Flash mikrokontrolera. Opracowano dedykowane funkcje w języku C/C++, których zadaniem jest odpakowanie tych struktur oraz przeprowadzenie wnioskowania opierając się na sprzętowych możliwościach mikroprocesora. Poniżej opisano zastosowane struktury danych oraz mechanikę działania poszczególnych algorytmów.

\subsection{Struktura i reprezentacja wag modeli implementacyjnych}
Podstawowym wyzwaniem wdrożenia modeli analitycznych na mikrokontrolerze ESP32 (dysponującym zaledwie 520 KB pamięci SRAM) jest konieczność wyeliminowania dynamicznej alokacji pamięci oraz struktur opartych na wskaźnikach, z których natywnie korzystają biblioteki w języku Python. Aby to osiągnąć, wyuczone na PC parametry modeli oraz współczynniki skalowania zostały wyekstrahowane i spłaszczone do postaci jednowymiarowych, statycznych tablic w języku C/C++. Użycie modyfikatora \texttt{const} sprawia, że wiedza modelu jest mapowana bezpośrednio do pamięci Flash mikrokontrolera (PROGMEM) redukując obciążenie pamięci operacyjnej podczas fazy wnioskowania. Poniżej opisano struktury wyeksportowane dla każdego z trzech algorytmów.

\subsubsection{ Autoenkoder}
W klasycznych bibliotekach sieci neuronowych takich jak scikit-learn, warstwy reprezentowane są jako obiekty operujące na dwuwymiarowych macierzach. W przygotowanym pliku nagłówkowym topologię modelu zredukowano do kilku  warstw neuronów. Wejście z 52 czujników kompresowane są sekwencyjnie do kolejnych warstw neuronów i rekonstruowane z powrotem do 52 wartości. Macierze wag zostały spłaszczone do jednowymiarowych tablic. Taki układ umożliwia mikrokontrolerowi wykonywanie propagacji sygnału w przód za pomocą prostych, zagnieżdżonych pętli \texttt{for}, bez alokacji macierzy pośrednich. Plik zawiera również wartość \texttt{THRESHOLD} wyznaczoną w oparciu o dystrybucję błędu średniokwadratowego (MSE) na zbiorze treningowym.

\begin{lstlisting}[language=C++, caption={Fragment pliku nagłówkowego z wagami modelu Autoenkodera.}]
#ifndef MODEL_WEIGHTS_H
#define MODEL_WEIGHTS_H

const int INPUT_DIM = 52;
const int NUM_LAYERS = 4;
const float THRESHOLD = 0.012959;

// Wektory do skalowania danych wejściowych
const float scaler_min[52] = {0.000000, 22.439240, 33.159720, 31.640620, 2.798032, /* ... */};
const float scaler_range[52] = {2.549016, 33.333330, 22.873270, 16.579870, 797.201968, /* ... */};

// Przykład definicji wag (macierze i wektory obciążeń warstw)
// const float w_layer_0[...] = { ... };
// const float b_layer_0[...] = { ... };

#endif
\end{lstlisting}

\subsubsection { Las Izolacyjny (Isolation Forest)}
Drzewa decyzyjne w bibliotece scikit-learn są w pamięci PC rozległymi grafami połączonymi wskaźnikami. W implementacji dla ESP32 zbiór 100 wyuczonych drzew lasu izolacyjnego został zserializowany do jednej dużej tablicy struktur \texttt{IsoNode}. Eliminacja wskaźników wymagała zastąpienia ich indeksami bezwzględnymi (\texttt{left}, \texttt{right}), określającymi dokładną pozycję dziecka w głównej tablicy. Zmienna \texttt{tree\_offsets} informuje mikrokontroler pod którym indeksem rozpoczyna się korzeń danego drzewa. Gdy wskaźniki \texttt{left} i \texttt{right} przyjmują wartość -2, ESP32 rozpoznaje liść i na podstawie zapamiętanej w nim ilości próbek (\texttt{samples}) estymuje długość niezbudowanej ścieżki algorytmem $c(n)$ \cite{isoforest}.

\begin{lstlisting}[language=C++, caption={Struktura danych i początkowe węzły modelu Isolation Forest.}]
#ifndef MODEL_WEIGHTS_H
#define MODEL_WEIGHTS_H

struct IsoNode {
    int feature;       // Indeks cechy podziału (-2 dla liścia)
    float threshold;   // Wartość progowa podziału
    int left;          // Indeks lewego syna
    int right;         // Indeks prawego syna
    int samples;       // Liczba próbek w węźle
};

const int INPUT_DIM = 52;
const int N_TREES = 100;
const int MAX_SAMPLES = 256;
const float THRESHOLD = 0.45f;

// Skalery zbieżne z architekturą przetwarzania sygnałów
const float scaler_min[52] = {0.000000, 22.439240, 33.159720, /* ... */};
const float scaler_range[52] = {2.549016, 33.333330, 22.873270, /* ... */};

// Indeksy startowe poszczególnych drzew w tablicy węzłów
const int tree_offsets[100] = {0, 143, 254, 381, 490, 625, /* ... */};

// Tablica przechowująca strukturę całego lasu (węzły wewnętrzne oraz liście)
const IsoNode iforest_nodes[10496] = {
    {1, 0.750228f, 1, 94, 256},
    {25, 0.815373f, 2, 39, 100},
    {41, 0.031913f, 3, 6, 31},
    {41, 0.012535f, 4, 5, 2},
    {-2, -2.000000f, -1, -1, 1}, // Przykładowy liść
    /* ... */
};

#endif
\end{lstlisting}

\subsubsection{ Drzewa Półprzestrzenne (Half-Space Trees)}
Model HS-Trees charakteryzuje się stałą strukturą drzew o zadanej wysokości (\texttt{TREE\_HEIGHT = 10}). Ze względu na pełną architekturę drzew binarnych, właściwości węzłów (takie jak indeksy testowanych cech, progi podziałów oraz masy/liczniki próbek w fazie referencyjnej) zostały rozdzielone na niezależne, zoptymalizowane pamięciowo tablice jednowymiarowe.

\subsubsection{ Drzewa Półprzestrzenne (Half-Space Trees)}
W modelu HS-Trees z biblioteki River zastosowano architekturę tzw. Struktury Tablic (ang. \textit{Struct of Arrays} - SoA). Zamiast używać pojedynczej tablicy przechowującej struktury (jak w Isolation Forest), parametry węzłów zostały rozbite na pięć równoległych, niezależnych wektorów (m.in. \texttt{node\_feature}, \texttt{node\_threshold}, \texttt{node\_r\_mass}). Taki zabieg nie tylko ściśle dopasowuje się do narzuconej, stałej głębokości drzew (\texttt{TREE\_HEIGHT = 10}), ale również optymalizuje współczynnik trafień do pamięci podręcznej (cache hit rate) mikroprocesora Xtensa w układzie ESP32. Ponadto w modelu tym nie estymuje się matematycznej głębokości cięcia. Kiedy podczas wnioskowania mikrokontroler dotrze do liścia, odczytuje bezpośrednio przypisaną mu tablicową wartość masy (\texttt{node\_r\_mass}), co stanowi kryterium gęstości profilowanego zachowania maszyny.

\begin{lstlisting}[language=C++, caption={Struktura płaska parametrów algorytmu HS-Trees.}]
#ifndef MODEL_WEIGHTS_H
#define MODEL_WEIGHTS_H

const int INPUT_DIM = 52;
const int NUM_TREES = 10;
const int TREE_HEIGHT = 10;
const float THRESHOLD = 0.65;

const float scaler_min[52] = {0.000000, 22.439240, 33.159720, /* ... */};
const float scaler_range[52] = {2.549016, 33.333330, 22.873270, /* ... */};

// Przesunięcia adresowe dla 10 zapisanych drzew binarnego podziału
const int tree_offsets[10] = {0, 2047, 4094, 6141, 8188, 10235, 12282, 14329, 16376, 18423};

// Rozbite tablice parametrów węzłów drzew struktur HS-Trees
const int node_feature[20470] = {33, 14, 38, 46, 21, 10, 1, 34, 11, 42, -2, -2, 42, /* ... */};
const float node_threshold[20470] = {0.167508, 0.306248, 0.623690, 0.210857, 0.170858, /* ... */};

#endif
\end{lstlisting}

\subsection{Implementacja algorytmów}
W tym podrozdziale opisano proces implementacji algorytmów detekcji anomalii bezpośrednio na mikrokontrolerze ESP32. Zaprezentowano w nim kody źródłowe oraz mechanizmy, które odtwarzają modele \textbf{wyuczone uprzednio w środowisku PC}, pozwalając na ich autonomiczne działanie w warunkach ograniczonej pamięci platformy ESP32.

\subsubsection{ Implementacja modelu Autoenkodera na ESP32}

Architektura wielowarstwowego perceptronu MLP w topologii autoenkodera została przeniesiona na mikrokontroler w postaci płaskich tablic reprezentujących macierze wag (WEIGHTS\_1, WEIGHTS\_2, etc.) oraz wektory obciążeń (BIAS\_1, BIAS\_2, etc.).

Proces detekcji rozpoczyna się od normalizacji sygnału wejściowego za pomocą estymatorów transformacji liniowej MinMaxScaler wyznaczonych w fazie treningu. Następnie, sygnał poddawany jest propagacji w przód (ang. forward pass) przez warstwy sieci neuronowej: warstwę wejściową, warstwy ukryte dokonujące kompresji reprezentacji oraz warstwę wyjściową. W każdym neuronie realizowana jest suma iloczynów wejść i wag, a wynik poddawany jest nieliniowej funkcji aktywacji ReLU. Ewaluacja anomalii opiera się na obliczeniu błędu średniokwadratowego (MSE) pomiędzy sygnałem wejściowym a zrekonstruowanym na wyjściu sieci. Wartość MSE wyższa od zdefiniowanego progu THRESHOLD klasyfikuje wektor danych jako stan anomalii maszyny.

\begin{lstlisting}[language=C++, caption={Odtworzenie algorytmu autoencodera na ESP32}]
#include "model_weights.h"

float input_data[INPUT_DIM];
float hidden_layer[HIDDEN_DIM];
float output_data[INPUT_DIM];

(...)

void loop() {
    if (Serial.available() > 0) {
        
        // Wczytywanie danych z portu szeregowego
        for (int i = 0; i < INPUT_DIM; i++) {
            input_data[i] = Serial.parseFloat();
        }
        while(Serial.read() != -1);

        // Skalowanie danych wejściowych (MinMaxScaler)
        for (int i = 0; i < INPUT_DIM; i++) {
            if (scaler_range[i] > 0) {
                input_data[i] = (input_data[i] - scaler_min[i]) / scaler_range[i];
            } else {
                input_data[i] = 0.0;
            }
        }

        // Propagacja w przód - Warstwa Ukryta
        for (int i = 0; i < HIDDEN_DIM; i++) {
            float sum = BIAS_1[i]; 
            for (int j = 0; j < INPUT_DIM; j++) {
                sum += input_data[j] * WEIGHTS_1[j * HIDDEN_DIM + i];
            }
            hidden_layer[i] = relu(sum);
        }

        // Propagacja w przód - Warstwa Wyjściowa
        for (int i = 0; i < INPUT_DIM; i++) {
            float sum = BIAS_2[i]; 
            for (int j = 0; j < HIDDEN_DIM; j++) {
                sum += hidden_layer[j] * WEIGHTS_2[j * INPUT_DIM + i];
            }
            output_data[i] = relu(sum);
        }

        // Obliczanie błędu rekonstrukcji (MSE)
        float mse = 0.0;
        for (int i = 0; i < INPUT_DIM; i++) {
            float diff = input_data[i] - output_data[i];
            mse += (diff * diff);
        }
        mse /= INPUT_DIM;

        (...)
        
        if (mse > THRESHOLD) {
            Serial.println(" -> WYNIK: ANOMALIA!");
        } else {
            Serial.println(" -> WYNIK: NORMALNY.");
        }
    }
}
\end{lstlisting}

\subsection{Wnioskowanie z wykorzystaniem algorytmu Isolation Forest}

Model Isolation Forest został zaadaptowany do środowiska ESP32 przez spłaszczenie struktur drzewiastych do jednowymiarowej tablicy struktur IsoNode. Każdy węzeł zawiera indeks analizowanej cechy, progową wartość podziału oraz wskaźniki na potomne gałęzie lewą i prawą. Tablica pomocnicza tree\_offsets pełni funkcję rejestru adresów korzeni poszczególnych drzew w przestrzeni pamięci stałej.

Algorytm procesuje znormalizowany wektor danych poprzez iteracyjne schodzenie po strukturze poszczególnych drzew. Dla każdego węzła realizowana jest instrukcja warunkowa decydująca o kierunku propagacji w dół grafu, z inkrementacją licznika głębokości. W przypadku napotkania węzła końcowego (liścia) zawierającego agregację próbek treningowych, algorytm aplikuje matematyczną korektę szacującą nierozwiniętą ścieżkę. Wyznaczony indeks anomalii (ang. Anomaly Score) obliczany jest na podstawie średniej głębokości ścieżek z całego zespołu drzew w odniesieniu do średniej teoretycznej dla zadanego rozmiaru podzbioru. Analogicznie do autoenkodera, przekroczenie ustalonego progu determinuje detekcję anomalii  \cite{isoforest}.

\begin{lstlisting}[language=C++, caption={Odtworzenie algorytmu Isolation Forest na ESP32}]
#include <Arduino.h>
#include <math.h>
#include "model_weights.h"

float input_data[INPUT_DIM];

// Funkcja c(n) - średnia długość ścieżki
float c_factor(int n) {
    if (n > 2) {
        return 2.0f * (log((float)(n - 1)) + 0.5772156649f) - (2.0f * (float)(n - 1) / (float)n);
    } else if (n == 2) {
        return 1.0f;
    }
    return 0.0f;
}

// Funkcja przechodzenia pojedynczego drzewa
float compute_path_length(const float* x, int tree_idx) {
    int node_idx = 0;
    int depth = 0;
    int offset = tree_offsets[tree_idx];
    
    while (true) {
        IsoNode node = iforest_nodes[offset + node_idx];
        
        // Cecha -2 : liście  (Scikit-Learn)
        if (node.feature == -2) {
            return (float)depth + c_factor((int)node.samples);
        }
        
        // Nawigacja w lewo lub w prawo w zależności od progu
        if (x[node.feature] <= node.threshold) {
            node_idx = node.left;
        } else {
            node_idx = node.right;
        }
        depth++;
    }
}

void loop() {

    if (Serial.available() > 0) {
        
    (...) 

        // --- START POMIARU CZASU ---
        unsigned long start_time = micros();

        // Skalowanie danych wejściowych
        for (int i = 0; i < INPUT_DIM; i++) {
            if (scaler_range[i] > 0) {
                input_data[i] = (input_data[i] - scaler_min[i]) / scaler_range[i];
            } else {
                input_data[i] = 0.0;
            }
        }

        // Obliczenie średniej długości ścieżki we wszystkich drzewach
        float avg_path = 0.0;
        for (int i = 0; i < N_TREES; i++) {
            avg_path += compute_path_length(input_data, i);
        }
        avg_path /= (float)N_TREES;

        // Obliczenie Anomaly Score: 2^(-E(h(x)) / c(n))
        float anomaly_score = pow(2.0f, -avg_path / c_factor(MAX_SAMPLES));

        // --- KONIEC POMIARU CZASU ---
        unsigned long end_time = micros();

        Serial.print("Anomaly Score: ");
        Serial.print(anomaly_score, 6);
        
        if (anomaly_score > THRESHOLD) {
            Serial.print(" | Status: !!! ANOMALIA !!!");
        } else {
            Serial.print(" | Status: OK");
        }

        (...)
    }
}
\end{lstlisting}

\subsection{Analiza z wykorzystaniem algorytmu Half-Space Trees}

Przeniesienie algorytmu HS-Trees na platformę mikrokontrolera zrealizowano z użyciem równoległych tablic stałych określających indeksy cech (node\_feature), progi podziałów półprzestrzeni (node\_threshold), relacje topologiczne (node\_left, node\_right) oraz masę węzłów (node\_r\_mass). W przeciwieństwie do poprzednich metod, HS-Trees ma predefiniowaną głębokość ustaloną dla każdego drzewa co ułatwia zarządzanie pamięcią.

Przepływ danych polega na trawersacji znormalizowanego sygnału przez drzewa, kierując się prostymi warunkami logicznymi odniesionymi do wartości progowych w węzłach. Zamiast mierzyć głębokość przejścia (Isolation Forest), zwiększana jest masa węzła (r\_mass), do którego przypisana została dana próbka. Masa odzwierciedla częstotliwość występowania podobnych wzorców w zbiorze. Niska wartość masy wskazuje, że próbka trafiła do rejonu przestrzeni cech rzadko eksplorowanego podczas fazy uczenia, co z wysokim prawdopodobieństwem świadczy o wystąpieniu anomalii. Ostateczna decyzja klasyfikacyjna zapada po agregacji wartości wag z zespołu wszystkich drzew.

\begin{lstlisting}[language=C++, caption={Odtworzenie algorytmu Half-space trees na ESP32}]
#include "model_weightsHS_Trees.h"

// Funkcja Min-Max Scaler
void scale_input(const float* raw_input, float* scaled_output) {
    for (int i = 0; i < INPUT_DIM; i++) {
        if (scaler_range[i] != 0) {
            scaled_output[i] = (raw_input[i] - scaler_min[i]) / scaler_range[i];
        } else {
            scaled_output[i] = 0;
        }
    }
}

// score dla pojedynczego drzewa
float score_tree(const float* input, int tree_idx) {
    int current_node_idx = tree_offsets[tree_idx];
    
    while (true) {
        int feature = node_feature[current_node_idx];
        
        // feature == -2: liść
        if (feature == -2) {
            return node_r_mass[current_node_idx];
        }

        // Decyzja o przejściu lewo/prawo
        if (input[feature] <= node_threshold[current_node_idx]) {
            int left_rel_offset = node_left[current_node_idx];
            if (left_rel_offset == -1) return node_r_mass[current_node_idx];
            current_node_idx += left_rel_offset;
        } else {
            int right_rel_offset = node_right[current_node_idx];
            if (right_rel_offset == -1) return node_r_mass[current_node_idx];
            current_node_idx += right_rel_offset;
        }
    }
}


float get_hs_trees_score(const float* scaled_input) {
    float total_r_mass = 0;
    
    // suma r_mass ze wszystkich drzew
    for (int i = 0; i < NUM_TREES; i++) {
        total_r_mass += score_tree(scaled_input, i);
    }
    
    // Normalizacja wyniku
    const float WINDOW_SIZE = 500.0;
    float avg_mass = total_r_mass / NUM_TREES;
    float normalized_score = 1.0 - (avg_mass / WINDOW_SIZE);
    
    return normalized_score;
}

(...)

void loop() {

    if (Serial.available() > 0) {
        String inputStr = Serial.readStringUntil('\n');
        float raw_sensor_data[INPUT_DIM];
        int valIdx = 0;
        
        // Parsowanie linii CSV ze złącza szeregowego
        char* ptr;
        char buf[inputStr.length() + 1];
        inputStr.toCharArray(buf, sizeof(buf));
        ptr = strtok(buf, ",");
        
        while (ptr != NULL && valIdx < INPUT_DIM) {
            raw_sensor_data[valIdx++] = atof(ptr);
            ptr = strtok(NULL, ",");
        }

        if (valIdx == INPUT_DIM) {
            float scaled_data[INPUT_DIM];
            unsigned long start_time = micros(); // Pomiar czasu predykcji
            
            scale_input(raw_sensor_data, scaled_data);
            float score = get_hs_trees_score(scaled_data);
            
            unsigned long end_time = micros(); 
            bool is_anomaly = (score > THRESHOLD);
            
            // Wypisanie wyniku
            Serial.print("Score: ");
            Serial.print(score, 6);
            if (is_anomaly) {
                Serial.print(" | Status: !!! ANOMALIA !!!");
            } else {
                Serial.print(" | Status: OK");
            }
            Serial.print(" | Time: "); Serial.print(end_time - start_time);
            Serial.print(" us");
            
        }
    }
}
\end{lstlisting}

\section{Wyniki detekcji i parametry wydajnościowe}

\vspace{0.5cm}

\begin{table}[htbp]
    \centering
    \caption{Metryki jakości klasyfikacji i macierz pomyłek}
    \vspace{0.2cm}
    \resizebox{\textwidth}{!}{%
    \begin{tabular}{@{} l l c c c c c r r r r @{}}
    \toprule
    \textbf{Platforma} & \textbf{Algorytm} & \textbf{Thresh.} & \textbf{Prec.} & \textbf{Rec.} & \textbf{F1} & \textbf{Acc.} & \textbf{TP} & \textbf{FP} & \textbf{TN} & \textbf{FN} \\
    \midrule
    PC & Autoencoder & 0.0081 & 0.0287 & 0.9737 & 0.0558 & 0.9621 & 74 & 2502 & 63518 & 2 \\
    PC & Isolation Forest & 0.0778 & 0.0753 & 0.7895 & 0.1375 & 0.9886 & 60 & 737 & 65283 & 16 \\
    PC & Half-Space Trees & 0.9762 & 0.0987 & 0.8684 & 0.1772 & 0.9907 & 66 & 603 & 65417 & 10 \\
    \midrule
    Raspberry Pi & Autoencoder & 0.0081 & 0.0287 & 0.9737 & 0.0558 & 0.9621 & 74 & 2502 & 63518 & 2 \\
    Raspberry Pi & Isolation Forest & 0.0769 & 0.0756 & 0.8158 & 0.1384 & 0.9883 & 62 & 758 & 65262 & 14 \\
    Raspberry Pi & Half-Space Trees & 0.9762 & 0.0987 & 0.8684 & 0.1772 & 0.9907 & 66 & 603 & 65417 & 10 \\
    \midrule
    Jetson & Autoencoder & 0.0081 & 0.0287 & 0.9737 & 0.0558 & 0.9621 & 74 & 2502 & 63518 & 2 \\
    Jetson & Isolation Forest & 0.0778 & 0.0753 & 0.7895 & 0.1375 & 0.9886 & 60 & 737 & 65283 & 16 \\
    Jetson & Half-Space Trees & 0.9762 & 0.0987 & 0.8684 & 0.1772 & 0.9989 & 66 & 603 & 65417 & 10 \\
    \midrule
    \multicolumn{11}{c}{\textit{ESP - Model nauczony na PC, przepuszczenie przez wagi (Inference)}} \\
    \midrule
    ESP & Autoencoder & 0.0081 & 0.0287 & 0.9737 & 0.0558 & 0.9621 & 74 & 2502 & 63518 & 2 \\
    ESP & Isolation Forest & 0.0778 & 0.0753 & 0.7895 & 0.1375 & 0.9886 & 60 & 737 & 65283 & 16 \\
    ESP & Half-Space Trees & 0.9762 & 0.0987 & 0.8684 & 0.1772 & 0.9907 & 66 & 603 & 65417 & 10 \\
    \bottomrule
    \end{tabular}%
    }
\end{table}

\vspace{0.5cm}

\begin{table}[htbp]
    \centering
    \caption{Parametry wydajnościowe i obciążenie zasobów}
    \vspace{0.2cm}
    \resizebox{\textwidth}{!}{%
    \begin{tabular}{@{} l l r r r r r r r r r @{}}
    \toprule
    \textbf{Platforma} & \textbf{Algorytm} & \textbf{Tren.} & \textbf{Detekcja} & \textbf{Lat. śr.} & \textbf{P95} & \textbf{P99} & \textbf{Przepust.} & \textbf{CPU} & \textbf{RAM} & \textbf{Temp.} \\
     & & \textbf{[s]} & \textbf{[s]} & \textbf{[ms]} & \textbf{[ms]} & \textbf{[ms]} & \textbf{[op/s]} & \textbf{[\%]} & \textbf{[MB]} & \textbf{[$^\circ$C]} \\
    \midrule
    PC & Autoencoder & 9.855 & 0.0454 & 0.0671 & 0.0785 & 0.1041 & 1456109 & 13.0 & 3.535 & brak \\
    PC & Isolation Forest & 1.319 & 0.1308 & 4.0783 & 5.257 & 5.853 & 505489 & 21.6 & 68.474 & brak \\
    PC & Half-Space Trees & 134.540 & 13.8561 & 0.2080 & 0.2301 & 0.2375 & 4770 & 9.8 & 268.981 & brak \\
    \midrule
    Raspberry Pi & Autoencoder & 18.223 & 0.1471 & 0.0844 & 0.0892 & 0.1002 & 449323 & 25.1 & 3.544 & 49.0 \\
    Raspberry Pi & Isolation Forest & 1.722 & 0.2364 & 6.2460 & 6.350 & 6.913 & 279590 & 25.5 & 44.001 & 46.9 \\
    Raspberry Pi & Half-Space Trees & 233.411 & 16.2614 & 0.2460 & 0.2680 & 0.2777 & 4064 & 23.2 & 268.981 & 47.4 \\
    \midrule
    Jetson & Autoencoder & 27.308 & 0.3307 & 0.2419 & brak & 0.2924 & 199888 & 58.7 & 3.483 & brak \\
    Jetson & Isolation Forest & 2.815 & 0.4014 & 10.9220 & brak & 11.925 & 164659 & 18.5 & 46.593 & brak \\
    Jetson & Half-Space Trees & 281.192 & 40.5304 & 0.6243 & brak & 0.7175 & 1630.78 & 17.1 & 400.220 & brak \\
    \midrule
    \multicolumn{11}{c}{\textbf{\textit{ESP - UWAGA: Model nauczony na PC, przepuszczenie przez wagi (wyłącznie inferencja)}}} \\
    \midrule
    ESP & Autoencoder & 9.855 & 0.000062 & 0.0873 & 0.1725 & 0.2009 & 11450 & brak & 0.0417 & 65.0 \\
    ESP & Isolation Forest & 1.319 & 0.001072 & 1.0792 & 1.4383 & 1.5156 & 926 & brak & 0.0417 & 64.8 \\
    ESP & Half-Space Trees & 134.540 & 0.000272 & 0.2320 & 0.3063 & 0.3245 & 4310 & brak & 0.0422 & 65.0 \\
    \bottomrule
    \end{tabular}%
    }
\end{table}

\section{Podsumowanie i wnioski końcowe}


\clearpage

\addcontentsline{toc}{section}{Literatura}

\begin{thebibliography}{99}
\bibitem{rpi5specs} \url{https://www.raspberrypi.com/documentation/computers/raspberry-pi.html}, Dostęp 02.06.2026.
\bibitem{jetsonspecs} \url{https://docs.nvidia.com/jetson/orin-nano-devkit/user-guide/latest/hardware_layout.html}, Dostęp 02.06.2026.
\bibitem{amdspecs} \url{https://www.amd.com/en/products/processors/laptop/ryzen/7000-series/amd-ryzen-7-7435hs.html}, Dostęp 02.06.2026.
\bibitem{esp32} \url{https://documentation.espressif.com/esp32\_datasheet\_en.html}, Dostęp 02.06.2026.
\bibitem{anomalysurvey} Chandola, Varun \& Banerjee, Arindam \& Kumar, Vipin. (2009). Anomaly Detection: A Survey. ACM Comput. Surv.. 41. 10.1145/1541880.1541882.
\bibitem{isoforest}Liu, Fei Tony \& Ting, Kai \& Zhou, Zhi-Hua. (2012). Isolation-Based Anomaly Detection. ACM Transactions on Knowledge Discovery From Data - TKDD. 6. 9-15. doi: 10.1145/2133360.2133363. 
\url{https://www.researchgate.net/publication/239761771\_Isolation-Based\_Anomaly\_Detection}, Dostęp 02.06.2026.
\bibitem{fastanomalydetection} Swee Chuan Tan, Kai Ming Ting, and Tony Fei Liu. 2011. Fast anomaly detection for streaming data. In Proceedings of the Twenty-Second international joint conference on Artificial Intelligence - Volume Volume Two (IJCAI'11). AAAI Press, 1511–1516.
\url{https://www.ijcai.org/Proceedings/11/Papers/254.pdf}, Dostęp 02.06.2026.
\bibitem{kaggledata} \url{https://www.kaggle.com/datasets/nphantawee/pump-sensor-data}, Dostęp 02.06.2026.
\bibitem{shi2016edge} Shi, W., Cao, J., Zhang, Q., Li, Y., \& Xu, L. (2016). Edge computing: Vision and challenges. \textit{IEEE Internet of Things Journal}, 3(5), 637-646.
\bibitem{capogrosso2024tinyml} Capogrosso, L., Cunico, F., Cheng, D. S., Fummi, F., \& Cristani, M. (2024). A Machine Learning-Oriented Survey on Tiny Machine Learning. \textit{IEEE Access}, 12, 23406-23426.
\bibitem{reis2025edge} Reis, M. J. C. S. (2025). Lightweight Signal Processing and Edge AI for Real-Time Anomaly Detection in IoT Sensor Networks. \textit{Sensors}, 25(21), 6629.
  
\end{thebibliography}

\clearpage

\end{document} 
